% =====================================================================
%  Subsection: Finite and infinite numbers of degrees of freedom
%  Paste inside the chapter, e.g. at the end of the section on Noether's
%  theorems.
%
%  Needs these macros (preamble of the chapter file):
%    \dd, \pd, \eps, \weakeq, \vx, \Stab, \Diff
%  and packages: amsmath, booktabs, enumitem.
%
%  It cross-references these labels (adapt or delete if absent):
%    sec:definitions, sec:noether, sec:galilei, sec:sr, sec:soldering,
%    sec:kepler, sec:law-solution, prop:main, ex:param, ex:noether2-em,
%    eq:noether1, eq:noether2, eq:param-identity
% =====================================================================
\subsection{Finite and infinite numbers of degrees of freedom}
\label{sec:dof}

The chapter moves freely between mechanical systems (the Galilei group, the
Kepler problem, the relativistic particle) and field theories
(electrodynamics, general relativity). This is legitimate because the two are
not different theories: a mechanical system is a field theory on a
one-dimensional base manifold, and a field theory is a mechanical system whose
coordinates carry a continuous label. This subsection makes both statements
explicit and notes where the reduction fails to be a pure special case.

\paragraph{Mechanics as a field theory on $M=\mathbb R_t$.}
Take the base manifold to be the time axis, with the single coordinate
$x^0\equiv t$, and identify
\begin{equation}\label{eq:dictionary-id}
  \varphi_a(x^\mu)\;\leftrightarrow\;q^i(t),\qquad
  \mathcal L(\varphi_a,\partial_\mu\varphi_a;x)\;\leftrightarrow\;L(q^i,\dot q^i,t),\qquad
  \partial_\mu\;\leftrightarrow\;\frac{\dd}{\dd t}.
\end{equation}
The formulas of Sections~\ref{sec:definitions} and~\ref{sec:noether} then
specialize as in Table~\ref{tab:dictionary-dof}.

\begin{table}[ht]
\centering
\small
\renewcommand{\arraystretch}{1.35}
\setlength{\tabcolsep}{4pt}
\begin{tabular}{@{}>{\raggedright\arraybackslash}p{7.3cm}>{\raggedright\arraybackslash}p{7.3cm}@{}}
\toprule
Field theory on $M^d$ & Mechanics ($d=1$)\\
\midrule
Fields $\varphi_a(x^\mu)$ & Coordinates $q^i(t)$\\
Euler--Lagrange: $\pd{\mathcal L}{\varphi_a}-\partial_\mu\pd{\mathcal L}{(\partial_\mu\varphi_a)}=0$
  & $\pd{L}{q^i}-\dfrac{\dd}{\dd t}\pd{L}{\dot q^i}=0$\\
Noether current: $J^\mu=\pd{\mathcal L}{(\partial_\mu\varphi_a)}X_a-K^\mu$,\ \ $\partial_\mu J^\mu\weakeq0$
  & Noether charge: $Q=p_iX^i-K$,\ \ $\dot Q\weakeq0$\\
Energy--momentum tensor: $T^\mu{}_\nu=\pd{\mathcal L}{(\partial_\mu\varphi_a)}\partial_\nu\varphi_a-\delta^\mu_\nu\mathcal L$
  & $T^0{}_0=p_i\dot q^i-L=H$\\
Explicit $x$-dependence: $\partial_\mu T^\mu{}_\nu\weakeq-\big(\pd{\mathcal L}{x^\nu}\big)_{\!\rm expl}$
  & $\dot H=-\pd{L}{t}$\\
\bottomrule
\end{tabular}
\caption{Mechanics as the case $d=1$ of field theory.}
\label{tab:dictionary-dof}
\end{table}

In $d=1$ the divergence $\partial_\mu J^\mu$ is a single derivative, so the
current \emph{is} the charge and the conservation law is constancy in time.
The derivation of Noether's first theorem is the same line in both cases. For
$\delta q^i=\eps X^i$ with $\delta L=\eps\dot K$,
\begin{equation}\label{eq:noether-mech}
  \delta L=\Big[\pd{L}{q^i}-\frac{\dd}{\dd t}\pd{L}{\dot q^i}\Big]\eps X^i
  +\eps\,\frac{\dd}{\dd t}\big(p_iX^i\big)=\eps\dot K
  \;\Longrightarrow\;
  \frac{\dd}{\dd t}\big(p_iX^i-K\big)\weakeq0,\qquad p_i=\pd{L}{\dot q^i},
\end{equation}
which is \eqref{eq:noether1} with one base coordinate. The Galilean boost of
Section~\ref{sec:galilei}, $X^i=u^it$ and $K=m\,\bm u\!\cdot\!\vx$, is an
application of exactly this formula.

\paragraph{The covariance machinery in $d=1$.}
Proposition~\ref{prop:main} applies with $\mathcal G=\Diff(\mathbb R)$, the
reparametrizations of time, and with the preferred chart, that is the one-form
$\dd t$, as the background $b_0$. A reparametrization $t\to f(t)$ pulls $\dd t$
back to $f'(t)\,\dd t$, so
\begin{equation}
  \Stab_{\Diff(\mathbb R)}(\dd t)=\{\,t\to t+c\,\}.
\end{equation}
The symmetry is time translation, and Noether's theorem gives conservation of
$H$. A prescribed time-dependent external force is a second background; it
shrinks the stabilizer to the identity and produces $\dot H=-\partial L/\partial t$,
the last row of Table~\ref{tab:dictionary-dof}. Noether's \emph{second}
theorem applies to the parametrized action of Example~\ref{ex:param} (or of
Section~\ref{sec:sr}), where $t$ is promoted to a dynamical variable: the
identity \eqref{eq:param-identity} is \eqref{eq:noether2} for
$\Diff(\mathbb R)$. The soldering form of Section~\ref{sec:soldering} reduces
to the single function $e_\lambda=\dd X^0/\dd\lambda=\dot t$, with
$D=\dot t^{-1}\dd/\dd\lambda$; the parametrized Newtonian action is
$\int\dot t\,\tfrac m2(D\vx)^2\,\dd\lambda$, which equals $\int\tfrac m2\dot\vx^2/\dot t\,\dd\lambda$.

\paragraph{Base symmetries and target symmetries.}
A mechanical system is a map $q\colon M^1\to N$ from the time axis into a
\emph{target} space $N$ (a sigma model). Transformations then split into those
acting on the base, which change the arguments of the fields, and those acting
on the target, which change their values. Table~\ref{tab:base-target} shows
that the distinction covers every example of the chapter.

\begin{table}[ht]
\centering
\footnotesize
\renewcommand{\arraystretch}{1.3}
\setlength{\tabcolsep}{3pt}
\begin{tabular}{@{}>{\raggedright\arraybackslash}p{3.0cm}%
                  >{\raggedright\arraybackslash}p{2.1cm}%
                  >{\raggedright\arraybackslash}p{2.4cm}%
                  >{\raggedright\arraybackslash}p{3.0cm}%
                  >{\raggedright\arraybackslash}p{3.6cm}@{}}
\toprule
System & Base $M$ & Target / fibre & Base symmetries & Target symmetries\\
\midrule
Newtonian particle & $\mathbb R_t$ & $E^3$ & time translation
  & $E(3)$; boosts are target maps with $t$-dependent parameter\\
Relativistic particle & $\mathbb R_\lambda$ & Minkowski space
  & $\Diff(\mathbb R)$ (gauge) & Poincar\'e\\
Real scalar field & Minkowski space & $\mathbb R$ & Poincar\'e
  & $\varphi\to-\varphi$ if $V$ is even\\
$O(N)$ scalar model & Minkowski space & $\mathbb R^N$ & Poincar\'e & $O(N)$\\
Gravity (tetrad form) & $M^4$ & fields $e^a{}_\mu$ & $\Diff(M)$ (gauge)
  & local Lorentz (gauge)\\
\bottomrule
\end{tabular}
\caption{Base and target symmetries.}
\label{tab:base-target}
\end{table}

\paragraph{Field theory as a mechanical system with a continuum of degrees of
freedom.}
In the opposite direction, a $3{+}1$ split of a field theory turns it into a
mechanical system: set $q_{\vx}(t)=\varphi(\vx,t)$, so that the discrete index
$i$ becomes the continuous label $\vx$. Sums become integrals,
\begin{equation}
  L=\int\dd^3x\,\mathcal L,\qquad
  \pi(\vx)=\frac{\delta L}{\delta\dot\varphi(\vx)},\qquad
  Q=\sum_ip_iX^i-K\ \longrightarrow\ \int\dd^3x\,\big(\pi X-K^0\big)=\int\dd^3x\,J^0 .
\end{equation}
A chain of masses shows how the two descriptions relate. On a ring of $N$
masses,
\begin{equation}
  L=\sum_{n}\Big[\tfrac12m\,\dot q_n^2-\tfrac12k\,(q_{n+1}-q_n)^2\Big].
\end{equation}
Set $q_n(t)=\varphi(na,t)$, $m=\rho a$, $k=\tau/a$, and keep the length $Na$
fixed. As $a\to0$,
\begin{equation}
  L\ \to\ \int\dd x\,\Big[\tfrac12\rho\,\dot\varphi^2-\tfrac12\tau\,\varphi'^2\Big],
  \qquad\rho\ddot\varphi=\tau\varphi'' .
\end{equation}
The same system has three different symmetry statements.
\begin{enumerate}[leftmargin=1.8em,itemsep=3pt,label=\textbf{\arabic*.}]
\item \emph{Shift} $q_n\to q_n+\eps$ for all $n$, a target symmetry with $X=1$,
  $K=0$. The charge $\sum_np_n=\sum_nm\dot q_n$ becomes $\int\rho\dot\varphi\,\dd x$,
  with the local current $J^t=\rho\dot\varphi$, $J^x=-\tau\varphi'$, and
  $\partial_tJ^t+\partial_xJ^x=\rho\ddot\varphi-\tau\varphi''=0$.
\item \emph{Lattice translation} $n\to n+1$, a discrete symmetry (the cyclic group
  $\mathbb Z_N$ of the ring). It has no Noether charge in the sense of
  Section~\ref{sec:noether}.
\item \emph{Continuum translation} $x\to x+b$, a base symmetry that exists only
  in the limit. By \eqref{eq:noether1} with $\delta\varphi=-b\,\partial_x\varphi$
  its charge is the field momentum $P=-\int\rho\,\dot\varphi\,\varphi'\,\dd x$.
\end{enumerate}
So the relation between the finite and the infinite system is not only
``special case.'' The finite chain has \emph{fewer} symmetries than its
continuum limit, because discretization generically breaks continuous base
symmetries (and, for a relativistic theory on a lattice, Lorentz invariance).

\paragraph{Where the reduction is not a pure special case.}
\begin{itemize}[leftmargin=1.6em,itemsep=2pt]
\item \emph{Boundary conditions.} For $d=1$ the only boundaries are the initial
  and final times, and $Q$ needs no fall-off condition. For $d>1$ the charge
  $Q=\int J^0\,\dd^3x$ is conserved only if the flux through spatial infinity
  vanishes, which requires conditions on the fields.
\item \emph{Constraints.} A gauge parameter $\eps(\vx,t)$ yields one constraint
  per spatial point (Gauss's law, $\nabla\!\cdot\vE=4\pi\rho$, in
  Example~\ref{ex:noether2-em}), whereas reparametrization in $d=1$ yields a
  single constraint.
\item \emph{Soldering form.} It has $d$ components ($d=4$ for spacetime), not one.
\item \emph{Manifest covariance.} A $3{+}1$ split gives up manifest covariance,
  and boosts then act nonlinearly, exactly as for the gauge-fixed particle
  $\lambda=t$ of Section~\ref{sec:sr}.
\item \emph{Symmetry breaking in the quantum theory.} With finitely many
  degrees of freedom the Stone--von Neumann theorem makes all irreducible
  representations of the canonical commutation relations unitarily
  equivalent, so a symmetry is implemented unitarily and a degenerate classical
  minimum is resolved by tunneling (the ground state of a symmetric double well
  is symmetric). With infinitely many degrees of freedom inequivalent
  representations exist, and spontaneous symmetry breaking of the quantum
  vacuum (Section~\ref{sec:kepler}) becomes possible. Classically, the
  distinction between symmetry of the law and symmetry of the solution
  (Section~\ref{sec:law-solution}) is already present in both cases.
\end{itemize}
In both directions Proposition~\ref{prop:main} is unchanged: the symmetry group
is the stabilizer of the background in the covariance group. Only the
dimension of the base and the supply of target symmetries differ.
